Transformers excel at \textit{in-context learning} (ICL)---learning from demonstrations without parameter updates---but how they do so remains a mystery. 
Recent work suggests that Transformers may internally run \gd\ (GD), a first-order optimization method, to perform ICL.
In this paper, we instead demonstrate that Transformers learn to approximate second-order optimization methods for ICL.
For in-context linear regression, Transformers share a similar convergence rate as \textit{\newton's Method}; both are exponentially faster than GD.
Empirically, predictions from successive Transformer layers closely match different iterations of Newton's Method \textit{linearly}, with each middle layer roughly computing 3 iterations; thus, Transformers and Newton's method converge at roughly the same rate. 
In contrast, \gd\ converges \textit{exponentially} more slowly.
We also show that Transformers can learn in-context on ill-conditioned data, a setting where \gd\ struggles but \newton\ succeeds. 
Finally, to corroborate our empirical findings, we prove that Transformers can implement $k$ iterations of Newton's method with $k + \mathcal{O}(1)$ layers.
\blfootnote{Our codes are available at \url{https://github.com/DeqingFu/transformers-icl-second-order}.}